Die Vektoren $\vec a$ und $\vec b$ spannen ein Parallelogramm auf. Wird dieses
entlang von $\vec c$ verschoben, entsteht ein Spat. Wir möchten nun das Volumen
dieses Spats berechnen.

\noindent
\begin{minipage}[c]{0.57\linewidth}
  \[
    \begin{aligned}
      V
      &=A h \\
      &=|\vec a\times\vec b|\,h\\
      &=|\vec a\times\vec b|\,|\vec c|\,|\cos(\varphi)|\\
      &=\bigl|(\vec a\times\vec b)\cdot\vec c\bigr|.
    \end{aligned}
  \]
\end{minipage}\hfill
\begin{minipage}[c]{0.39\linewidth}
  \centering
  \CachedTikzFigure{content/resources/chapters/vector_algebra/tikz/scalar_triple_product_volume.tex}
\end{minipage}

Motiviert durch diese Überlegung, definieren wir das Spatprodukt.
\begin{DefinitionBox}{Spatprodukt}{vector_algebra_scalar_triple_product}
  Das \emph{Spatprodukt} \EnglishTerm{scalar triple product} der Vektoren
  $\vec a,\vec b,\vec c\in\R^3$ ist der Skalar
  \[
    S=(\vec a\times\vec b)\cdot\vec c.
  \]
  Der Betrag des Spatprodukts, $|S|$, entspricht dem Volumen des aufgespannten
  Spats.
  \[
    V = |S| = |(\vec a\times\vec b)\cdot\vec c|.
  \]
\end{DefinitionBox}

\begin{ExampleBox}{Volumen eines Spats}{vector_algebra_scalar_triple_product_volume_example}
  Die Vektoren $\vec a, \vec b, \vec c \in \R^3$ spannen einen Spat auf.
  \[
    \vec a=\begin{pmatrix}3\\-1\\2\end{pmatrix},\qquad
    \vec b=\begin{pmatrix}4\\2\\1\end{pmatrix},\qquad
    \vec c=\begin{pmatrix}0\\2\\-5\end{pmatrix}
  \]
  Das Spatprodukt $S$ ergibt sich zu
  \[
    S=(\vec a\times\vec b)\cdot\vec c
    =\begin{ClozeContinuation}[29]{3}
      \left(
        \begin{pmatrix}3\\-1\\2\end{pmatrix}
        \times
        \begin{pmatrix}4\\2\\1\end{pmatrix}
      \right)
      \cdot\begin{pmatrix}0\\2\\-5\end{pmatrix}
      =\begin{pmatrix}-5\\5\\10\end{pmatrix}
       \cdot\begin{pmatrix}0\\2\\-5\end{pmatrix}
      =-40
    \end{ClozeContinuation}.
  \]
  Damit beträgt das Volumen
  \[
    V=\Cloze{9cm}{|S|=|-40|=40}.
  \]
\end{ExampleBox}

\clearpage
Analog zu Kreuzprodukt lässt sich auch das Spatprodukt in Determinantenschreibweise
darstellen. Wir werden im Folgenden zeigen, dass
Mit der Determinantenschreibweise des Kreuzprodukts aus
Merke~\ref{key:vector_algebra_cross_product_determinant_notation} rechnen wir
die Determinantendarstellung des Spatprodukts explizit nach. Wegen
$\vec a\cdot\vec e_i=a_i$ folgt durch Entwicklung nach der ersten Spalte
\begin{ClozeArea}[34]{13}
  \[
    \begin{aligned}
      (\vec a\times\vec b)\cdot\vec c
      &=\vec a\cdot(\vec b\times\vec c)\\[1.2em]
      &=a_1\det\!\begin{pmatrix}b_2&c_2\\b_3&c_3\end{pmatrix}
        -a_2\det\!\begin{pmatrix}b_1&c_1\\b_3&c_3\end{pmatrix}
        +a_3\det\!\begin{pmatrix}b_1&c_1\\b_2&c_2\end{pmatrix}
        \\[1.2em]
      &=a_1(b_2c_3-b_3c_2)
        -a_2(b_1c_3-b_3c_1)
        +a_3(b_1c_2-b_2c_1)
        \\[1.2em]
      &=\det\!\begin{pmatrix}
          a_1&b_1&c_1\\
          a_2&b_2&c_2\\
          a_3&b_3&c_3
        \end{pmatrix}
        =\det(\vec a,\vec b,\vec c).
    \end{aligned}
  \]
\end{ClozeArea}

\begin{KeyPointBox}{Spatprodukt als Determinante}{vector_algebra_scalar_triple_product_determinant}
  Für
  \[
    \vec a=\begin{pmatrix}a_1\\a_2\\a_3\end{pmatrix},\qquad
    \vec b=\begin{pmatrix}b_1\\b_2\\b_3\end{pmatrix},\qquad
    \vec c=\begin{pmatrix}c_1\\c_2\\c_3\end{pmatrix}
  \]
  gilt
  \[
    (\vec a\times\vec b)\cdot\vec c
    =\begin{ClozeContinuation}[20]{3}
      \det(\vec a,\vec b,\vec c)
      =
      \det\!\begin{pmatrix}
        a_1&b_1&c_1\\
        a_2&b_2&c_2\\
        a_3&b_3&c_3
      \end{pmatrix}
    \end{ClozeContinuation}.
  \]
  Die Vektoren $\vec a$, $\vec b$ und $\vec c$ bilden dabei die Spalten der
  Matrix. Der Betrag der Determinante beschreibt das Volumen des aufgespannten
  Spats:
  \[
    V=\Cloze{12cm}{
      \bigl|(\vec a\times\vec b)\cdot\vec c\bigr|
      =\bigl|\det(\vec a,\vec b,\vec c)\bigr|
    }.
  \]
\end{KeyPointBox}
